\documentclass[11pt,a4paper]{article}

% ===========================================================================
% DeltaField -- Paper 2 (the "danach" paper).
% STATUS: SKELETON. Method, theorem and protocol are written from
% docs/DELTAFIELD_CONCEPT.md and the shipped code (v5_rankbind, head='deltafield',
% v5_rankbind/tests/test_deltafield.py). The RESULTS section is a deliberate
% placeholder: only one single-seed BRENDA-200 pilot has finished
% (matrix MRR 0.206, matrix per-ligand AUC 0.926, anti-shortcut holds) and it
% ranks below the v4 recipe, so no headline result is claimed yet. The
% GNN-deepened difference propagation, the multi-seed runs and the enzyme-wide
% (turnover) runs are the contribution still being produced; fill the result
% tables once those land.
% ===========================================================================

% ---- Packages -------------------------------------------------------------
\usepackage[utf8]{inputenc}
\usepackage[T1]{fontenc}
\usepackage{microtype}
\usepackage[margin=1in]{geometry}
\usepackage{amsmath,amssymb,amsfonts}
\usepackage{graphicx}
\usepackage{booktabs}
\usepackage{multirow}
\usepackage{array}
\usepackage{xcolor}
\usepackage{authblk}
\usepackage{caption}
\usepackage{subcaption}
\usepackage{enumitem}
\usepackage{url}
\usepackage[hidelinks,colorlinks=true,linkcolor=black,citecolor=blue,urlcolor=blue]{hyperref}
\usepackage[capitalize,noabbrev]{cleveref}
\usepackage{natbib}
\usepackage{lmodern}
\usepackage{xspace}
\usepackage{amsthm}

\graphicspath{{../figures/}{figures/}}

% ---- Macros ---------------------------------------------------------------
\newcommand{\df}{DeltaField\xspace}
\newcommand{\method}{RankBind\xspace}
\newcommand{\nullprior}{\textsc{null\_prot\_prior}\xspace}
\newcommand{\gauc}{global AUC}
\newcommand{\plauc}{per-ligand AUC}
\newcommand{\matmrr}{matrix MRR}
\newcommand{\Hfree}{\mathbf{H}_{\text{free}}}
\newcommand{\Hcoup}{\mathbf{H}_{\text{coup}}}

% "good"/"bad" cells + a visible placeholder marker for unfilled results.
\newcommand{\good}[1]{\textcolor{teal!75!black}{\textbf{#1}}}
\newcommand{\bad}[1]{\textcolor{red!75!black}{\textbf{#1}}}
\newcommand{\TODO}[1]{\textcolor{red!80!black}{\textbf{[TODO: #1]}}}
\newcommand{\pending}{\textcolor{gray}{---}}

% ---- Submission placeholders ----------------------------------------------
\newcommand{\repourl}{\url{https://github.com/cmanitz/rankbind}}
\newcommand{\zenododoi}{10.5281/zenodo.XXXXXXX}

\newtheorem{theorem}{Theorem}
\newtheorem{proposition}{Proposition}

% ---- Title block ----------------------------------------------------------
\title{\textbf{\df: Architecturally Ligand-Conditional Binding Prediction\\
via Coupled--Free Difference Encoding}\\[4pt]
\large\normalfont\textcolor{red!70!black}{Draft skeleton --- results pending}}

\author[1]{Christoph Manitz}
\affil[1]{Leipzig University, Germany \\
\texttt{christoph.manitz@uni-leipzig.de}}

\date{June 2026}

% ===========================================================================
\begin{document}
\maketitle

% ---- Abstract -------------------------------------------------------------
\begin{abstract}
Drug--target interaction models on enzyme--substrate benchmarks pass pooled
discrimination metrics (global AUC) by exploiting a \emph{protein-level
shortcut}: they infer the binding probability mostly from the protein
identity, with the ligand a secondary signal. Companion
work~\citep{manitz2026rankbind} diagnoses this shortcut and shows that a
four-ingredient \emph{training recipe} (\method) breaks it, at the cost of
giving up the pooled metric. That fix is enforced by the loss and the
sampler; nothing in the architecture \emph{prevents} the score from
re-encoding a per-protein prior. We propose \df, an architecture in which the
protein-only prior is \emph{structurally unrepresentable in the score}.
A single weight-shared bipartite residue$\times$atom transformer is run twice
on the same inputs --- once with all cross-modal edges masked (the
\emph{free} pass, which can encode only ligand-agnostic protein identity) and
once with them open (the \emph{coupled} pass) --- and the score is read
\emph{only} from the per-node difference field $\Delta = \Hcoup - \Hfree$.
Any signal invariant to the ligand's presence lives in the free pass and
cancels in the subtraction; the surviving signal is ligand-conditional by the
chain rule rather than by a penalty. This yields a verifiable guarantee
(\Cref{thm:zerofield}): with the cross-edges masked, $\Delta\equiv0$ exactly
and the score collapses to a single global constant identical for every
protein, so a per-protein prior \emph{cannot} rank --- a property we assert as
a unit test rather than a claim. \df\ drops into the existing \method\ stack
(margin loss, hard-negative mining, $200\times200$ matrix-ranking evaluation)
and additionally emits a residue/atom/contact interaction map read off the
same difference field. This manuscript presents the architecture, the
zero-field theorem and the evaluation protocol; the full empirical study
---a GNN-deepened difference-propagation variant, multi-seed runs, and the
enzyme-wide transfer--- is in progress and the result tables are placeholders.
A single-seed BRENDA-200 pilot already confirms the anti-shortcut property
(matrix per-ligand AUC $0.926$) but does not yet match the \method\ recipe on
raw ranking; closing that gap is the object of the ongoing work.
\end{abstract}

% ===========================================================================
\section{Introduction}
\label{sec:intro}

Predicting whether a small molecule binds a given protein is a central task
in computational chemistry. On the enzyme--substrate sub-task, several
architectures report pooled AUC values that suggest saturated
benchmarks~\citep{bai2023drugban,huang2021moltrans,nguyen2021graphdta,wang2023gems}.
Companion work~\citep{manitz2026rankbind} shows these numbers certify a
\emph{shortcut}~\citep{geirhos2020shortcut}: a data-blind classifier that
scores each pair using only the per-protein training positive rate
($\nullprior$) reproduces both the pooled AUC and the attractor geometry
(Gini $\approx0.995$) of the trained baselines, while every baseline's
\emph{within-ligand} ranking sits at or below chance. The pathology is a
\emph{rank} pathology that pooled metrics cannot see.

\paragraph{From a training fix to an architectural one.}
\method~\citep{manitz2026rankbind} removes the shortcut with four training
ingredients --- protein-balanced sampling, a within-ligand margin loss, a
low-rank bilinear head, and online hard-negative mining --- and demonstrably
drives the top-10 attractor overlap with $\nullprior$ to zero. But that
guarantee is \emph{behavioural}: it is a property of a converged optimum
under a particular loss and sampler. The bilinear score
$s(L,P)=f(L)^{\top}(UV^{\top}+\mathrm{diag}(d))g(P)+b$ still contains a
protein-only pathway --- $g(P)$ enters the score whether or not the ligand
carries information --- so a per-protein prior remains \emph{representable};
the recipe merely makes it unattractive to the optimiser. On a new dataset,
a new encoder, or a tuning slip, the prior can re-enter (we observe exactly
this on one external corpus in the companion paper, where training
early-stops into a shortcut-flavoured solution).

This paper asks a sharper question: \emph{can the protein-only prior be made
unrepresentable in the score by construction?} A score that is identically a
constant whenever the ligand is removed cannot encode a per-protein prior at
all --- no loss or sampler required. We give such an architecture.

\paragraph{The \df\ idea.}
We borrow the difference-encoding device of GMSF~\citep{zhao2026gmsf}, which
predicts enzyme function from a chemical reaction by encoding the
\emph{per-atom difference} between aligned reactant and product states: the
difference localises the reaction centre and cancels the conserved scaffold.
The transferable kernel is \emph{represent a process by the per-element
difference between two aligned states; the difference localises the active
region and cancels the inert background}. Binding has two analogous states
--- unbound (apo) and bound (holo) --- but no holo structure exists at
inference. Rather than hallucinate a virtual complex, \df\ \emph{realises both
states inside one forward pass} of a single weight-shared network, toggling
only the cross-modal attention mask:
\begin{itemize}[leftmargin=*,itemsep=1pt,topsep=2pt]
  \item the \textbf{free pass} masks all ligand$\leftrightarrow$protein edges,
        so its representation can encode at most ligand-agnostic protein
        identity --- exactly the shortcut;
  \item the \textbf{coupled pass} opens those edges with the same weights;
  \item the score is read from the \textbf{difference field}
        $\Delta=\Hcoup-\Hfree$, in which the shortcut cancels.
\end{itemize}

\paragraph{Contributions.}
\begin{enumerate}[leftmargin=*,itemsep=0pt,topsep=2pt]
  \item \textbf{\df}, a weight-shared bipartite residue$\times$atom
        transformer whose score is read \emph{only} from the coupled--free
        difference field, making the per-protein prior structurally
        unrepresentable in the score (\Cref{sec:method}).
  \item A \textbf{verifiable anti-shortcut guarantee}: the zero-field theorem
        (\Cref{thm:zerofield}) and its corollary that the score collapses to
        a single global constant under masked coupling, asserted as a unit
        test (\Cref{sec:antishortcut}).
  \item A \textbf{dual output}: a residue/atom/contact interaction map read
        off the same difference field, with a pre-registered pocket-overlap
        go/no-go that tests whether the field concentrates \emph{on} the
        binding pocket --- the site that \method's learned attention was
        shown to \emph{avoid}~\citep{manitz2026rankbind} (\Cref{sec:dual}).
  \item An \textbf{evaluation protocol} that reuses the companion paper's
        diagnosis instruments ($\nullprior$ probe, $200\times200$ matrix
        ranking, matched-capacity ablation) so \df\ is measured on the same
        axes as the recipe it is meant to internalise (\Cref{sec:exp}). The
        empirical study is in progress; \Cref{sec:exp} states the
        pre-registered checks and leaves the result tables as placeholders.
\end{enumerate}

% ===========================================================================
\section{Related Work}
\label{sec:related}

\paragraph{Target-prior / shortcut bias in DTI.}
Shortcut learning~\citep{geirhos2020shortcut} and benchmark leakage in
protein--ligand prediction~\citep{wallach2018mostligands,kapoor2023leakage}
are well documented. The companion paper~\citep{manitz2026rankbind} measures
the specific protein-prior shortcut on enzyme--substrate data and removes it
behaviourally. \textbf{TAPB}~\citep{wang2025tapb} addresses the same
target-prior bias but removes it by \emph{causal backdoor adjustment} over a
dictionary of confounders. \df\ differs in mechanism: it removes the prior by
\emph{representation-level subtraction} (the free pass), not by an explicit
causal estimand.

\paragraph{Apo$\to$holo and counterfactual complexes.}
\textbf{FlowDock}~\citep{morehead2025flowdock} shares the apo$\to$holo framing
and an affinity head, but predicts on a \emph{generated} 3D complex. \df's
coupled pass is a counterfactual virtual-holo computed by \emph{masking}, not
a hallucinated structure, and its score is read from a difference in
embedding space with no 3D generation.

\paragraph{Residue$\times$atom interaction maps.}
\textbf{MONN}~\citep{li2020monn}, \textbf{LaMPSite}, \textbf{ICAN} and
\textbf{ArkDTA} produce residue$\times$atom cross-attention interaction maps.
\df\ borrows the map but reads its \emph{score} only from the difference
field, subject to the zero-field invariant; the cited models score from the
coupled representation directly, so they carry no analogous structural
anti-shortcut guarantee.

\paragraph{Difference encoding.}
The structural seed is \textbf{GMSF}~\citep{zhao2026gmsf}, which encodes the
per-atom reactant$\to$product difference and propagates it with a second GNN.
GMSF's own ablation is instructive: the difference module is the main
structural win, and \emph{naive} fusion of two unaligned modalities causes
negative transfer --- the alignment anchor is what makes fusion help. \df\
inherits both lessons (the difference as the load-bearing device; an explicit
alignment, here the residue/atom identity across the free and coupled passes).

% ===========================================================================
\section{\df: Method}
\label{sec:method}

\df\ reuses the frozen encoders and projectors of the \method\ stack and adds
a single new interaction head (\texttt{head\_type=\textquotesingle
deltafield\textquotesingle}); the training loop, loss and metrics are
unchanged. Tensor dimensions follow the repository: ESM2 per-residue
$1280$-d, ChemBERTa per-token $384$-d, projection $d$.

\subsection{Inputs and projection}
The protein (apo) is the frozen ESM2 per-residue tensor
$E_P\in\mathbb{R}^{B\times L\times1280}$ with a residue mask, as already
produced in \method's \texttt{attn\_pool} mode. The ligand is the frozen
ChemBERTa \emph{per-token} tensor $E_L\in\mathbb{R}^{B\times A\times384}$ with
a token mask (the per-token hidden states, minus the CLS token, cached
alongside a BPE-token$\to$RDKit-atom offset map). The existing projectors are
applied along the token/residue axis,
$u=\mathrm{ProteinProjector}(E_P)\in\mathbb{R}^{B\times L\times d}$,
$v=\mathrm{LigandProjector}(E_L)\in\mathbb{R}^{B\times A\times d}$; with
modality-type embeddings added, they are stacked into a single node set
$X_0=[\,u\,;\,v\,]\in\mathbb{R}^{B\times(L+A)\times d}$.

\subsection{Weight-shared bipartite encoder: free vs.\ coupled}
The engine is a small transformer of $T$ blocks (multi-head attention $+$
FFN), run \emph{twice with identical weights}, differing only in the
attention mask:
\begin{align}
\Hfree &= \mathrm{Enc}_\theta(X_0;\, A_{\text{free}}), &
\Hcoup &= \mathrm{Enc}_\theta(X_0;\, A_{\text{coup}}),
\end{align}
where $A_{\text{free}}$ is block-diagonal (residues attend to residues, atoms
to atoms; \emph{no} cross edges) and $A_{\text{coup}}$ additionally opens the
ligand$\leftrightarrow$protein edges. The free pass therefore cannot transport
any ligand information into a protein node, and vice versa: it can encode at
most ligand-agnostic protein identity --- the shortcut. The coupled
cross-attention weights $W_{\text{cross}}\in\mathbb{R}^{B\times \text{heads}
\times L\times A}$ are cached for the contact map.

\subsection{The difference field}
The load-bearing object is the per-node difference
\begin{equation}
\Delta \;=\; \Hcoup - \Hfree \;\in\; \mathbb{R}^{B\times(L+A)\times d},
\end{equation}
split into a residue block $d_{\text{res}}\in\mathbb{R}^{B\times L\times d}$
and an atom block $d_{\text{atom}}\in\mathbb{R}^{B\times A\times d}$, with
scalar per-node energies
$e_{\text{res}}=\lVert d_{\text{res}}\rVert_2$ and
$e_{\text{atom}}=\lVert d_{\text{atom}}\rVert_2$.

\subsection{Difference-graph propagation (the GNN graft)}
\label{sec:prop}
Following GMSF's ``second GNN'', an ablatable propagation stage
($\texttt{prop\_layers}\in\{0,1,\dots\}$) lifts the difference from a per-node
quantity to a synergistic, pocket-scale signal:
$x'_{\text{node}}=\mathrm{Linear}([\Hfree;\Delta])$ followed by one (or more)
weight-shared message-passing layer under the coupled mask, yielding a
propagated field $\tilde{D}$. \textbf{Only $\tilde{D}$ routes into the score;
$\Hfree$ is a context channel for the propagation only and never reaches the
scalar.} Deepening this stage --- replacing the single coupled-mask layer with
a genuine residue$\times$atom message-passing network over the contact graph
--- is the principal architectural extension under study
(\Cref{sec:future}); the zero-field guarantee below holds for any
$\tilde{D}$ that is a function of $\Delta$ alone and vanishes with it.

\subsection{Additive perturbation-energy score head}
The score is a mass-conserving sum of per-node contributions, every term
linear in $\tilde{D}$:
\begin{align}
w_{\text{res}} &= \mathrm{softmax}_i\!\big(\phi\cdot \tilde{D}_{\text{res}}\big)\ \text{gated by}\ e_{\text{res}}, &
c_{\text{res}}[i] &= a^{\top}\tilde{D}_{\text{res}}[i], \\
\text{score} &= b_0 \;+\; \sum_i w_{\text{res}}[i]\,c_{\text{res}}[i] \;+\; \sum_a w_{\text{atom}}[a]\,c_{\text{atom}}[a], \label{eq:score}
\end{align}
with an analogous atom branch. Crucially, \textbf{no $a^{\top}\Hfree$ or
$g(P)$ term ever enters \Cref{eq:score}}: the protein representation reaches
the score only through $\tilde{D}$, which is a difference. The head outputs
(a) the scalar score for matrix-MRR ranking, (b) the residue interaction map
$e_{\text{res}}$ / signed $c_{\text{res}}$, (c) the atom map $e_{\text{atom}}$
back-projected to SMILES atoms, and (d) the residue$\times$atom contact map
$C[i,a]=\mathrm{mean}_{\text{heads}}W_{\text{cross}}[i,a]\cdot
e_{\text{res}}[i]\cdot e_{\text{atom}}[a]$, matching \method's
\texttt{head(fL,gP)}$\to$scalar contract so the rest of the pipeline is
untouched.

\subsection{Cost and caching}
The free pass depends only on the protein, so $\Hfree(P)$ is cached per
protein and reused across all ligands; only the coupled pass and the field run
per (ligand, protein) chunk. This amortises both the $200\times200$
evaluation matrix and the per-epoch hard-negative refresh, which would
otherwise be prohibitive given the $L\times A$ cross-attention.

% ===========================================================================
\section{Anti-shortcut by construction}
\label{sec:antishortcut}

\df's anti-shortcut property is structural, not behavioural. Three locks,
each tied to a \emph{measured} failure mode in the companion
paper~\citep{manitz2026rankbind}:

\begin{enumerate}[leftmargin=*,itemsep=2pt,topsep=2pt]
  \item \textbf{The free pass is the shortcut, and it is subtracted out.}
        $\Hfree$ encodes only ligand-agnostic protein identity --- the exact
        signal that gives $\nullprior$ its Gini $\approx0.995$ and the
        Phase-1 baselines their top-10 attractor Jaccard $0.54$--$0.67$.
        $\Delta=\Hcoup-\Hfree$ cancels it.
  \item \textbf{No additive protein term reaches the scalar.} Remove the
        ligand and $\tilde{D}\to0$, so $\text{score}\to b_0$, a single global
        constant. A per-protein prior is \emph{literally unrepresentable} in
        the score --- the formal analogue of GMSF's conserved-scaffold
        cancellation.
  \item \textbf{Limited degrees of freedom.} A sparsity penalty, a
        total-variation locality penalty along the backbone, and the
        inherited within-ligand margin cap the protein's ability to smear a
        global identity across residues to fake conditionality; hard
        negatives (the same hot protein against competing ligands) force the
        field to \emph{differ across ligands} or the margin is violated.
\end{enumerate}

\begin{theorem}[Zero-field invariant]
\label{thm:zerofield}
Let the coupled mask equal the free mask (all cross-modal edges removed).
Then $\Hcoup=\Hfree$, hence $\Delta\equiv0$ and $\tilde{D}\equiv0$, and by
\Cref{eq:score} $\,\mathrm{score}=b_0$ for every protein in the batch.
Consequently $\mathrm{std}(\mathrm{score})=0$: the model cannot rank proteins,
so no per-protein prior is expressible in the score.
\end{theorem}

\begin{proof}[Proof sketch]
With identical masks the two encoder runs are the same function of the same
input, so $\Hcoup=\Hfree$ exactly and $\Delta=0$. Every term of
\Cref{eq:score} other than $b_0$ is linear in $\tilde{D}$, a function of
$\Delta$ that vanishes with it; therefore the sum is $b_0$ for every node
configuration. $b_0$ is a shared scalar, independent of $P$.
\end{proof}

\paragraph{The theorem is a unit test, not a claim.}
\Cref{thm:zerofield} is asserted in continuous integration
(\texttt{v5\_rankbind/tests/test\_deltafield.py::test\_zero\_field\_invariant}):
with coupling forced off, the test checks $e_{\text{res}}=e_{\text{atom}}=0$ to
$10^{-6}$, that the score equals $b_0$ for every protein, and that the score
standard deviation across the batch is below $10^{-6}$. Companion tests assert
the converse (coupling produces a non-zero field and ranking-capable scores),
ligand-conditionality (same protein, different ligand $\Rightarrow$ different
score), gradient flow, and matched parameter budget ($d=128\Rightarrow
\approx\!650$k trainable parameters, within $15\%$ of \method's $627{,}201$,
so ``more capacity'' is never the explanation). The anti-shortcut guarantee is
thus \emph{verifiable on every commit} rather than argued in prose.

% ===========================================================================
\section{Dual output and the interpretability go/no-go}
\label{sec:dual}

The interaction map is produced, not explained post-hoc: the residue map
$e_{\text{res}}$ (and signed $c_{\text{res}}$, a mass-conserving share of the
predicted affinity), the atom map $e_{\text{atom}}$, and the contact map
$C[i,a]$ all fall out of the same difference field that produces the score.
They are validated by feeding $e_{\text{res}}$ into the companion paper's
existing pocket-overlap audit
(\texttt{evaluation/attn\_annotation\_scan.py}): the within-protein ROC-AUC of
a per-residue score against the UniProt active-/binding-site mask.

\paragraph{Pre-registered go/no-go.}
The companion paper established that \method's learned attention
\emph{de-weights} the catalytic pocket (within-protein ROC-AUC $0.08$--$0.21$,
i.e.\ far below the chance value $0.5$): the learned weights track
hydrophobicity, not function. \df's difference field is built to put signal
\emph{on} the interface. We pre-register the go/no-go that $e_{\text{res}}$
achieve within-protein ROC-AUC $>0.5$ against the annotated pocket (ideally
$>0.7$), \emph{flipping} the $0.08$--$0.21$ result. This is reported as a
\textbf{secondary} metric: the ranking claim does not depend on it, and a null
result here (the field tracking biophysics rather than the annotated pocket,
as ESM2 is sequence-trained) is a documented risk (\Cref{sec:limitations}),
not a refutation of the anti-shortcut guarantee.

% ===========================================================================
\section{Experiments}
\label{sec:exp}

\paragraph{Protocol (fixed; results pending).}
\df\ is trained and evaluated on exactly the companion paper's instruments so
the two are directly comparable: the BRENDA protein-stratified split (seed
$42$); the $200\times200$ matrix-ranking metrics (matrix MRR, Hit@$K$); the
$\nullprior$ probe (top-10 Jaccard, per-row Spearman); matched capacity
($d=128$, $\approx\!650$k parameters vs.\ \method's $627{,}201$); the
within-ligand margin loss ($k=4$, $m=1.0$) on hard-negative triplets; three
seeds $\{42,7,1337\}$; AdamW / cosine / bf16; frozen PLMs. Field regularisers
(\Cref{sec:recipe}) are small and ablatable.

\paragraph{Pre-registered pass/fail checks.}
\begin{enumerate}[leftmargin=*,itemsep=1pt,topsep=2pt]
  \item[(A)] \textbf{Ranking.} Matrix MRR $\ge$ the \method\ v4 default
        ($0.326\pm0.072$) at matched capacity --- the difference-only score
        must not cost ranking.
  \item[(B)] \textbf{Anti-shortcut.} Top-10 Jaccard vs.\ $\nullprior$
        $\le 0.11$ (ideally $\approx0$) \emph{and} the zero-field unit test
        passes exactly.
  \item[(C)] \textbf{Interpretability (secondary).} $e_{\text{res}}$
        binding/active-site within-protein ROC-AUC $>0.5$, flipping the
        \method\ $0.08$--$0.21$.
\end{enumerate}

\paragraph{Status.}
A single-seed BRENDA-200 pilot (\texttt{tag v6\_deltafield}, run dir
\texttt{results/v5\_rankbind/20260610-121458\_*\_abl\_deltafield\_v6\_deltafield})
has finished and confirms check~(B) --- the matrix per-ligand AUC is
$0.926$, on par with \method\ v5b, and the architecture rides no prior --- but
\emph{not yet} check~(A): its matrix MRR is $0.206$, below the v4 recipe's
$0.326$. The difference-only score is therefore anti-shortcut as designed but
not yet ranking-competitive at this depth. The remaining runs --- the
GNN-deepened difference-propagation variant (\Cref{sec:prop}), the multi-seed
sweep, and the enzyme-wide (turnover) transfer --- are in progress; their
results populate the tables below.

\begin{table}[h]
  \centering
  \footnotesize
  \setlength{\tabcolsep}{6pt}
  \caption{\df\ vs.\ the \method\ recipe on BRENDA-200 (matched capacity,
  three seeds). \textbf{Placeholder} --- the multi-seed and GNN-deepened rows
  are pending; the single-seed pilot is shown for orientation only and is
  \emph{not} a headline result.}
  \label{tab:headline}
  \begin{tabular}{lcccc}
    \toprule
    Config & matrix MRR & Hit@10 & matrix \plauc & Top-10 Jac.\ vs.\ \nullprior \\
    \midrule
    \method\ v4 (reference)            & $0.326\pm0.072$ & $0.755\pm0.095$ & $0.891$ & $0.035\pm0.030$ \\
    \method\ v5b (reference)           & $0.427\pm0.123$ & $0.814\pm0.103$ & $0.930$ & $0.000$ \\
    \midrule
    \df\ (single-seed pilot, $s{=}42$) & $0.206$         & $0.559$         & $0.926$ & \TODO{read} \\
    \df\ ($\texttt{prop\_layers}{=}0$, 3 seeds) & \pending & \pending & \pending & \pending \\
    \df\ (GNN-deepened, 3 seeds)       & \pending & \pending & \pending & \pending \\
    \bottomrule
  \end{tabular}
\end{table}

\begin{table}[h]
  \centering
  \footnotesize
  \setlength{\tabcolsep}{6pt}
  \caption{Enzyme-wide transfer (BRENDA+SABIO \texttt{turnover}), matching the
  companion paper's transfer probe. \textbf{Placeholder} --- runs in flight.}
  \label{tab:transfer}
  \begin{tabular}{lcccc}
    \toprule
    Dataset / Config & matrix MRR & $\rho_{\text{row}}$ vs.\ \nullprior & Top-10 Jac. & zero-field test \\
    \midrule
    \texttt{turnover} \df\ (baseline) & \pending & \pending & \pending & \pending \\
    \texttt{turnover} \df\ (hard-neg) & \pending & \pending & \pending & \pending \\
    \bottomrule
  \end{tabular}
\end{table}

\paragraph{Interpretability result (secondary).}
\TODO{pocket-overlap ROC-AUC of $e_{\text{res}}$ vs.\ the v5b $0.08$--$0.21$;
report whether the go/no-go of \Cref{sec:dual} is met.}

% ===========================================================================
\section{Training recipe}
\label{sec:recipe}

\df\ slots into the \method\ v4 training loop. The primary loss is unchanged:
the within-ligand \texttt{margin\_loss} ($k=4$, $m=1.0$) on
\texttt{TripletCollator} triplets with hard-negative mining
(\texttt{hard\_pool\_size} scaled to the train-protein count, per the
companion paper's pool-coverage lesson), refreshed every $2$--$3$ epochs (the
field forward is heavier than a bilinear dot). Small, ablatable field
regularisers shape the difference field:
\begin{itemize}[leftmargin=*,itemsep=0pt,topsep=2pt]
  \item $\mathcal{L}_{\text{sparse}}=\lambda_s(\overline{e_{\text{res}}}+\overline{e_{\text{atom}}})$
        --- localise the interface (with $\lambda$ annealing and a floor on
        total field energy to avoid collapse);
  \item $\mathcal{L}_{\text{tv}}=\lambda_{\text{tv}}\,\mathrm{TV}(e_{\text{res}})$
        along the backbone --- a contiguous pocket, not speckle;
  \item $\mathcal{L}_{\text{neg}}=\lambda_{\text{neg}}\,
        \overline{\lVert\tilde{D}(L,P^-)\rVert^2}$ over hard negatives --- a
        non-binder must induce a near-zero field, training the score-collapse
        property directly;
  \item coupling-edge dropout ($p\approx0.15$) --- the GMSF AAM-noise
        transplant.
\end{itemize}
An optional, pretraining-only branch (off the critical path) adds focal-BCE on
the contact map $C[i,a]$ and distils $\mathrm{ESM2}(\text{holo})-
\mathrm{ESM2}(\text{apo})$ into $\tilde{D}$ \emph{if} PDBbind contacts become
available; \df\ must pass margin-only on BRENDA without any of this, since no
apo/holo pairs exist locally.

% ===========================================================================
\section{Limitations and open risks}
\label{sec:limitations}

\begin{itemize}[leftmargin=*,itemsep=0pt,topsep=1pt]
  \item \textbf{Ranking not yet competitive.} The single-seed pilot ranks
        below the \method\ recipe (matrix MRR $0.206$ vs.\ $0.326$); the
        anti-shortcut property holds but the difference-only score must be
        made to rank as well as the bilinear recipe. This is the central open
        problem and the motivation for the GNN-deepened propagation.
  \item \textbf{Sparsity collapse} (most likely failure mode): too-strong
        $\mathcal{L}_{\text{sparse}}/\mathcal{L}_{\text{neg}}$ drives
        $\Delta\to0$ everywhere and MRR to chance. Mitigated by $\lambda$
        annealing and a field-energy floor; monitored in
        \texttt{train\_log.jsonl}.
  \item \textbf{Free-pass leakage} via the $[\Hfree;\Delta]$ re-injection in
        the propagation layer --- if $\Hfree$ reaches the scalar the prior
        re-enters. Only $\tilde{D}$ routes to the score; gated on the Jaccard
        audit; the channel is ablated if Jaccard creeps up.
  \item \textbf{Inherited pocket anti-correlation.} ESM2 is sequence-trained;
        the field may track biophysics (hydrophobicity) rather than the
        interface even after subtraction (cf.\ the \method\ v5b finding).
        Pocket overlap is therefore a \emph{gating} experiment, not a headline.
  \item \textbf{Compute.} The coupled pass plus $L\times A$ cross-attention
        over $B\cdot k$ negatives is far heavier than a bilinear dot;
        free-pass caching per protein and entmax/top-$k$ truncation are
        load-bearing, and single-head attention with gradient checkpointing
        may be required.
  \item \textbf{ChemBERTa BPE $\ne$ heavy atoms.} The token$\to$atom map is
        approximate (subword-mean fallback); the atom-level map is framed as
        coarse and validated against RDKit substructures.
  \item \textbf{Novelty defence.} \df\ is novel \emph{as a composition} (the
        design panel's prior-art checks found no DTI/PLI paper with this
        device); the constituent pieces are borrowed and cited. The paper must
        lead with the concrete null-baseline failure mode and the CI-tested
        zero-field theorem, not the abstract ``shortcut'' label, or a reviewer
        will read it as derivative of MONN/TAPB.
\end{itemize}

% ===========================================================================
\section{Conclusion and planned work}
\label{sec:future}

\df\ makes the protein-level shortcut --- diagnosed and behaviourally removed
in companion work~\citep{manitz2026rankbind} --- \emph{structurally
unrepresentable} in the score, by reading the score only from the coupled--free
difference field of a single weight-shared bipartite transformer. The
guarantee is a verifiable theorem rather than a converged behaviour. This
manuscript fixes the architecture, the theorem and the protocol; the empirical
contribution is the object of ongoing work, in priority order:
\begin{enumerate}[leftmargin=*,itemsep=1pt,topsep=2pt]
  \item \textbf{Deepen the difference propagation (\Cref{sec:prop}).} Replace
        the single coupled-mask layer with a genuine residue$\times$atom
        graph network over the contact graph, to close the ranking gap to the
        \method\ recipe while preserving the zero-field invariant (which holds
        for any $\tilde{D}=\tilde{D}(\Delta)$).
  \item \textbf{Multi-seed BRENDA-200} ($\{42,7,1337\}$) plus the
        matched-capacity head ablation (\df\ vs.\ bilinear).
  \item \textbf{Enzyme-wide transfer} on BRENDA+SABIO \texttt{turnover}
        (runs in flight), to test whether the architectural guarantee
        survives the $50\times$ scale-up where the recipe needed pool-size
        rescaling.
  \item \textbf{Interpretability go/no-go} (\Cref{sec:dual}): does the
        difference field land on the pocket the recipe's attention avoided?
\end{enumerate}

% ===========================================================================
\section*{Declarations}

\paragraph{Data and software availability.}
\df\ is implemented in the same repository as the companion
work~\citep{manitz2026rankbind} (\texttt{head\_type=\textquotesingle
deltafield\textquotesingle} in \texttt{v5\_rankbind/model.py}, config
\texttt{v5\_rankbind/configs/abl\_deltafield.json}, CI in
\texttt{v5\_rankbind/tests/test\_deltafield.py}), released at \repourl\ and
archived under Zenodo DOI \texttt{\zenododoi}. All datasets, splits and
encoders are those of the companion paper.

\paragraph{Competing interests.} The author declares none.

\paragraph{Author contributions.} C.M.\ conceived and implemented the method
and wrote the manuscript.

% ===========================================================================
\bibliographystyle{plainnat}
\bibliography{../references,references_deltafield}

\end{document}
